% Generated by publication/build_sections.py; do not edit by hand.
\paragraph{Informal verdict.}
\textbf{A --- Complete solution}, with \textbf{high confidence} under the response's explicit convention that a solution is a \(C^2\) function on an open interval and that \(yy''=f\) means \(y(x)y''(x)=f(x)\).

\paragraph{Lean coverage.}
Full canonical coverage under the explicit classical C2 pointwise convention. For part \(a\), Lean proves that the explicit inverse-integral family gives entire solutions of y y\_second = 1 and that every interval solution is its restriction, hence exactly the maximal solutions. For the open-ended part \(b\), Lean proves forcing continuity and zero compatibility, simple and tangent zero jet constraints, the fundamental integral identity, the full f = 0 classification, and a smooth nonuniqueness example. The answer's additional local sufficiency and endpoint-asymptotic theorems are optional strengthenings, not requirements of the printed question.

\subsection{Audited informal answer}
Q830
\begin{enumerate}
\item Faithful restatement and interpretation
\end{enumerate}
The PDF asks:

\(a\) Determine the maximal solutions of \(y\,y''=1\).
\(b\) Given a map \(f:\mathbb R\to\mathbb R\), what can be said about the solutions of \(y\,y''=f\)?

There is no further hypothesis, footnote, or contextual qualification attached to Q830. The printed equation in part \(b\) is \(y\,y''=f\), not \(y\,y''=f(y)\); I therefore read it pointwise as

\[
y(x)y''(x)=f(x).
\]

Part \(b\) genuinely leaves unspecified the domain and regularity of a ``solution,'' whether global or local solutions are sought, and what kind of classification is expected. Luc Pommeret+1
I use the following natural convention.
A classical solution on an open interval \(I\) is a function \(y\in C^2(I,\mathbb R)\) satisfying the equation at every point of \(I\). A solution is maximal if it is not the restriction of a classical solution to a strictly larger open interval.

\begin{enumerate}
\item Part \(a\): complete classification of the maximal solutions
\end{enumerate}
Define

\[
\Psi(q):=\sqrt2\int_0^q e^{s^2}\,ds
=\sqrt{\frac{\pi}{2}}\operatorname{erfi}(q),
\qquad q\in\mathbb R.
\]

The function \(\Psi\) is odd, \(C^\infty\), strictly increasing, and maps \(\mathbb R\) bijectively onto \(\mathbb R\).
Theorem 1
The maximal classical solutions of

\[
y\,y''=1
\]

are precisely the functions parametrized by

\[
\boxed{
x=\tau+a\Psi(q),\qquad
y=\varepsilon a e^{q^2},
\qquad q\in\mathbb R,
}
\]

where

\[
a>0,\qquad \tau\in\mathbb R,\qquad
\varepsilon\in\{-1,1\}.
\]

Equivalently,

\[
\boxed{
y(x)=\varepsilon a
\exp\!\left(
\Psi^{-1}\!\left(\frac{x-\tau}{a}\right)^2
\right),
\qquad x\in\mathbb R.
}
\]

Every maximal solution is therefore defined on all of \(\mathbb R\). Its sign is constant, and \(|y|\) has a unique minimum \(a\), attained at \(x=\tau\).
Proof
Let \(y\in C^2(I)\) solve \(yy''=1\) on an interval \(I\).
Since the left-hand side would vanish wherever \(y=0\), one has

\[
y(x)\ne0\qquad(x\in I).
\]

Thus \(y\) has a constant sign. Put

\[
\varepsilon=\operatorname{sgn}y,\qquad u=\varepsilon y=|y|>0.
\]

Because \(y=\varepsilon u\) and \(y''=\varepsilon u''\),

\[
u\,u''=1.
\]

Multiplication by \(u'/u\), or direct differentiation, gives the first integral

\[
\frac{d}{dx}\left(\frac12u'^2-\log u\right)
=u'\left(u''-\frac1u\right)=0.
\]

Hence

\[
\frac12u'^2-\log u=C
\]

for some constant \(C\). Set

\[
a=e^{-C}>0.
\]

Then

\[
u'^2=2\log\frac ua.
\]

In particular,

\[
u\ge a.
\]

Define

\[
q=\frac{u'}{\sqrt2}.
\]

The first integral becomes

\[
u=a e^{q^2}.
\]

Moreover,

\[
q'=\frac{u''}{\sqrt2}
=\frac1{\sqrt2\,u}.
\]

Consequently,

\[
\frac{d}{dx}\bigl(a\Psi(q)\bigr)
=a\sqrt2\,e^{q^2}q'
=a\sqrt2\,\frac ua\frac1{\sqrt2u}
=1.
\]

Therefore there exists \(\tau\in\mathbb R\) such that

\[
x=\tau+a\Psi(q).
\]

Since \(y=\varepsilon u\), this gives

\[
y=\varepsilon a e^{q^2}.
\]

Conversely, let

\[
x=\tau+a\Psi(q),\qquad y=\varepsilon a e^{q^2}.
\]

Because

\[
\frac{dx}{dq}=a\sqrt2\,e^{q^2}>0,
\]

\(q\) is a smooth function of \(x\). Differentiating,

\[
y'
=
\frac{2\varepsilon aqe^{q^2}}
{a\sqrt2e^{q^2}}
=\varepsilon\sqrt2\,q,
\]

and hence

\[
y''
=\varepsilon\sqrt2\,\frac{dq}{dx}
=\frac{\varepsilon}{ae^{q^2}}
=\frac1y.
\]

Thus \(yy''=1\).
Finally, \(\Psi(q)\to\pm\infty\) as \(q\to\pm\infty\), so the parametrization gives a solution on all of \(\mathbb R\). Every local solution obtained above is a restriction of one of these entire solutions. Therefore these and only these are maximal. \(\square\)
Initial-value parametrization
For arbitrary data

\[
y(x_0)=y_0\ne0,\qquad y'(x_0)=v_0,
\]

the parameters are

\[
\varepsilon=\operatorname{sgn}y_0,\qquad
q_0=\frac{\varepsilon v_0}{\sqrt2},
\]


\[
a=|y_0|e^{-v_0^2/2},
\qquad
\tau=x_0-a\Psi(q_0).
\]

Thus every admissible initial condition determines exactly one maximal solution.

\begin{enumerate}
\item Part \(b\): general classical-solution theory
\end{enumerate}
Because \(f\) is allowed to be completely arbitrary, no universal explicit formula can exist. Indeed, for every prescribed \(Y\in C^2(I)\), choosing

\[
f=YY''
\]

makes \(Y\) a solution. Thus arbitrary growth, oscillation, zero sets, and flat behavior can be encoded in \(f\).
There is nevertheless a precise structural theory. It consists of:


ordinary existence and uniqueness wherever \(y\ne0\);


exact compatibility conditions at zeros;


an integral identity controlling consecutive zeros;


a sharp classification of finite singular endpoints when \(f\) has a nonzero limit.



3.1 Immediate necessary conditions
Let \(y\in C^2(I)\) satisfy

\[
yy''=f.
\]

Then necessarily

\[
\boxed{f\in C(I)}
\]

because \(f=yy''\) is a product of continuous functions.
Moreover,

\[
\boxed{y(c)=0\quad\Longrightarrow\quad f(c)=0.}
\]

Thus

\[
Z(y)\subseteq Z(f).
\]

In particular, if \(f\) is discontinuous on every nonempty open interval---for example \(f=\mathbf 1_{\mathbb Q}\)---then there is no classical solution on any nonempty open interval.

3.2 Regular branches: existence, uniqueness, and continuation
Assume now that \(f\in C(I)\), where \(I\) is an open interval.
Theorem 2
Given

\[
x_0\in I,\qquad y_0\ne0,\qquad v_0\in\mathbb R,
\]

there exists a unique maximal zero-free branch

\[
y:J\longrightarrow\mathbb R,
\qquad x_0\in J\subset I,
\]

such that

\[
y(x_0)=y_0,\qquad y'(x_0)=v_0,
\qquad yy''=f.
\]

On \(J\), it satisfies the Volterra equation

\[
\boxed{
y(x)=y_0+v_0(x-x_0)
+\int_{x_0}^x (x-t)\frac{f(t)}{y(t)}\,dt.
}
\]

If a finite endpoint \(c\in I\) of \(J\) exists, then

\[
\boxed{\liminf_{x\to c,\ x\in J}|y(x)|=0.}
\]

Proof
Away from \(y=0\), the equation is the first-order system

\[
\begin{cases}
y'=v,\\[2mm]
v'=\dfrac{f(x)}y.
\end{cases}
\]

Near \((x_0,y_0,v_0)\), choose a rectangle on which

\[
|y|\ge \frac{|y_0|}{2}.
\]

On this rectangle the vector field

\[
F(x,y,v)=\left(v,\frac{f(x)}y\right)
\]

is continuous and locally Lipschitz in \((y,v)\). The usual contraction argument applied to

\[
Y(x)=Y_0+\int_{x_0}^x F(t,Y(t))\,dt
\]

gives a unique local solution. Uniqueness allows all local continuations that remain in \(y\ne0\) to be united into one maximal zero-free branch.
Integrating \(y''=f/y\) twice gives the displayed Volterra equation.
Now suppose that \(c\in I\) is a finite endpoint and that

\[
\liminf_{x\to c}|y(x)|>0.
\]

Then \(|y|\ge\delta>0\) near \(c\). Since \(f\) is bounded near \(c\),

\[
|y''|\le \frac{\sup|f|}{\delta}.
\]

Hence \(y'\) has a finite limit at \(c\), and then \(y\) also has a finite nonzero limit. The local existence theorem at \(c\) extends the branch past \(c\), contradicting maximality. \(\square\)
A zero-free branch need not have a unique maximal classical continuation through a zero; this distinction is essential and will be illustrated below.

3.3 A fundamental integral identity
For every classical solution,

\[
(yy')'=y'^2+yy''=y'^2+f.
\]

Hence, for \(a<b\) in the solution interval,

\[
\boxed{
\int_a^b f(x)\,dx
=
y(b)y'(b)-y(a)y'(a)
-\int_a^b y'(x)^2\,dx.
}
\]

Consequences
If \(y(a)=y(b)=0\), then

\[
\boxed{
\int_a^b f(x)\,dx
=-\int_a^b y'(x)^2\,dx\le0.
}
\]

Equality occurs exactly when \(y\equiv0\) on \([a,b]\), in which case \(f\equiv0\) there.
Thus a nonzero solution cannot have two zeros \(a<b\) unless

\[
\int_a^b f<0.
\]

In particular, if \(f\ge0\), a solution cannot have two distinct zeros unless it vanishes identically, together with \(f\), between those zeros.
If \(y\) is \(T\)-periodic, then

\[
\int_x^{x+T} f(t)\,dt
=-\int_x^{x+T}y'(t)^2\,dt\le0.
\]


\begin{enumerate}
\item Exact local analysis at a zero
\end{enumerate}
Let \(c\) be a zero of a classical solution, and write

\[
p=y'(c),\qquad q=y''(c).
\]

Taylor's formula with Peano remainder gives

\[
y(c+h)
=ph+\frac12qh^2+o(h^2).
\]

4.1 Necessary jet conditions
Simple zero: \(p\ne0\)
Since \(y''(c+h)\to q\),

\[
\frac{f(c+h)}h
=
\frac{y(c+h)}h\,y''(c+h)
\longrightarrow pq.
\]

Therefore \(f\) is differentiable at \(c\), and

\[
\boxed{f'(c)=y'(c)y''(c)=pq.}
\]

Tangential zero: \(p=0\)
Then

\[
\frac{f(c+h)}{h^2}
=
\frac{y(c+h)}{h^2}\,y''(c+h)
\longrightarrow \frac12q^2.
\]

Hence

\[
\boxed{
f(c+h)=\frac12q^2h^2+o(h^2).
}
\]

If \(q\ne0\), then \(f\) is positive on both sides of \(c\), sufficiently close to \(c\).
If \(q=0\), necessarily

\[
\boxed{f(c+h)=o(h^2).}
\]

If zeros of \(y\) accumulate at \(c\), then necessarily

\[
y'(c)=y''(c)=0.
\]

Indeed, a simple or nondegenerate quadratic zero is isolated.

4.2 Complete characterization of simple zeros
Theorem 3
Let \(f\) be defined near \(c\). There exists a local classical solution with

\[
y(c)=0,\qquad y'(c)=p\ne0
\]

if and only if

\[
\boxed{
g(h):=\frac{f(c+h)}h
}
\]

extends continuously to \(h=0\).
When this condition holds, for every prescribed \(p\ne0\) there is a unique such local solution, and

\[
\boxed{
y''(c)=\frac{g(0)}p.
}
\]

Proof
Necessity was proved above.
For sufficiency, translate \(c\) to \(0\), so that

\[
f(h)=h\,g(h)
\]

with \(g\) continuous. Seek

\[
y(h)=h\,z(h),\qquad z(0)=p.
\]

The equation is equivalent to

\[
y''(h)=\frac{g(h)}{z(h)}.
\]

Twice integrating from \(0\), using \(y(0)=0\) and \(y'(0)=p\), gives

\[
h z(h)
=
ph+\int_0^h(h-s)\frac{g(s)}{z(s)}\,ds,
\]

or equivalently

\[
\boxed{
z(h)
=
p+h\int_0^1(1-r)\frac{g(hr)}{z(hr)}\,dr.
}
\]

Choose \(\delta>0\) and consider the complete metric space

\[
\mathcal B
=
\left\{
z\in C([-\delta,\delta]):
\|z-p\|_\infty\le\frac{|p|}{2}
\right\}.
\]

Every \(z\in\mathcal B\) satisfies \(|z|\ge |p|/2\).
Define

\[
(Tz)(h)
=
p+h\int_0^1(1-r)\frac{g(hr)}{z(hr)}\,dr.
\]

If \(M=\sup_{|h|\le\delta}|g(h)|\), then

\[
\|Tz-p\|_\infty\le\frac{\delta M}{|p|}.
\]

For sufficiently small \(\delta\), \(T\) maps \(\mathcal B\) into itself. Moreover,

\[
\|Tz_1-Tz_2\|_\infty
\le
\frac{2\delta M}{p^2}
\|z_1-z_2\|_\infty.
\]

Again decreasing \(\delta\), this is a strict contraction. Hence \(T\) has a unique fixed point \(z\).
Set \(y(h)=hz(h)\). The integral equation becomes

\[
y(h)
=
ph+\int_0^h(h-s)\frac{g(s)}{z(s)}\,ds.
\]

The integrand is continuous, so \(y\in C^2\) and

\[
y''(h)=\frac{g(h)}{z(h)}.
\]

Therefore

\[
y(h)y''(h)=hz(h)\frac{g(h)}{z(h)}
=hg(h)=f(h).
\]

Finally,

\[
y''(0)=\frac{g(0)}p.
\]

Uniqueness follows from uniqueness of the fixed point. \(\square\)
Thus simple zeros are completely understood.

4.3 Complete characterization of nondegenerate tangent zeros
Theorem 4
Let \(f\) be defined near \(c\). There is a local classical solution satisfying

\[
y(c)=y'(c)=0,\qquad y''(c)\ne0
\]

if and only if

\[
\boxed{
\frac{f(c+h)}{h^2}
}
\]

extends continuously to \(h=0\) with a strictly positive value

\[
A>0.
\]

When this holds, there are exactly two local solutions with a nondegenerate tangent zero at \(c\). They satisfy respectively

\[
\boxed{
y''(c)=+\sqrt{2A}
\qquad\text{and}\qquad
y''(c)=-\sqrt{2A}.
}
\]

Proof
Necessity follows from

\[
\frac{f(c+h)}{h^2}\longrightarrow\frac12y''(c)^2.
\]

For sufficiency and uniqueness, translate \(c\) to \(0\) and write

\[
f(h)=h^2g(h),\qquad g(0)=A>0.
\]

Choose

\[
B=\pm\sqrt{\frac A2}.
\]

We seek a solution of the form

\[
y(h)=h^2z(h),\qquad z(h)\longrightarrow B.
\]

It suffices first to work on \(h>0\). Put

\[
h=e^{-t},\qquad Z(t)=z(e^{-t}),\qquad G(t)=g(e^{-t}).
\]

Then \(t\to+\infty\) corresponds to \(h\to0^+\). A direct calculation gives

\[
y''=Z''-3Z'+2Z,
\]

where primes now denote \(t\)-derivatives. The equation becomes

\[
Z\bigl(Z''-3Z'+2Z\bigr)=G(t).
\]

Write

\[
Z=B+w.
\]

Because \(A=2B^2\), the equation is

\[
w''-3w'+4w
=
H(t,w),
\]

where

\[
H(t,\xi)
=
\frac{G(t)}{B+\xi}-2B+2\xi.
\]

As \(t\to\infty\),

\[
H(t,0)\to0,
\qquad
\partial_\xi H(t,0)
=
-\frac{G(t)}{B^2}+2\to0.
\]

Let

\[
\omega=\frac{\sqrt7}{2},
\qquad
K(r)=\frac1\omega e^{-3r/2}\sin(\omega r),
\qquad r\ge0.
\]

For a continuous function \(R\) tending to \(0\), the unique solution tending to \(0\) at \(+\infty\) of

\[
w''-3w'+4w=R
\]

is

\[
w(t)=\int_0^\infty K(r)R(t+r)\,dr.
\]

Consider the Banach space \(C_0([T,\infty))\) of continuous functions tending to \(0\), with the supremum norm, and define

\[
(\mathcal Tw)(t)
=
\int_0^\infty K(r)H(t+r,w(t+r))\,dr.
\]

Since \(H(t,0)\to0\) and its derivative with respect to its second variable tends uniformly to \(0\) near \(\xi=0\), one may choose \(T\) large and a small ball

\[
\|w\|_\infty\le\rho<\frac{|B|}{2}
\]

such that \(\mathcal T\) maps this ball into itself and is a strict contraction. Therefore there is a unique \(w\in C_0([T,\infty))\) solving the equation.
Returning to \(h=e^{-t}\) gives a unique right-hand solution with

\[
\frac{y(h)}{h^2}\to B.
\]

The same argument applied to \(f(-h)\) gives the unique left-hand solution. They glue in \(C^2\) at \(0\), because

\[
y(0)=y'(0)=0,\qquad
y''(0)=2B=\pm\sqrt{2A}.
\]

Conversely, any solution with this second derivative gives \(Z(t)\to B\) and \(Z'(t)\to0\); variation of constants forces it to satisfy the same fixed-point equation, so uniqueness follows from the contraction. \(\square\)
For example, when \(f(x)=x^2\), the two tangent solutions are

\[
y(x)=\pm\frac{x^2}{\sqrt2}.
\]


4.4 Degenerate zeros
The remaining case is

\[
y(c)=y'(c)=y''(c)=0.
\]

Necessarily

\[
f(c+h)=o(h^2),
\]

but no finite-order jet condition can classify all such solutions.
For instance, with

\[
\phi(x)=
\begin{cases}
e^{-1/x^2},&x\ne0,\\
0,&x=0,
\end{cases}
\]

the function \(\phi\) is \(C^\infty\) and flat at \(0\). Taking

\[
f=\phi\phi'',
\]

both \(\phi\) and various sign-gluings of \(\phi\) provide solutions with a completely flat zero. Thus the degenerate case genuinely contains arbitrary smooth flat behavior.

\begin{enumerate}
\item Finite endpoints of zero-free branches
\end{enumerate}
Let \(y\) be a maximal zero-free branch on an interval ending at a finite point \(c\), and put

\[
u=|y|>0.
\]

Then

\[
u\,u''=f.
\]

5.1 Positive limiting forcing repels the solution from zero
Theorem 5
If

\[
f(c)>0,
\]

then \(c\) cannot be a finite endpoint of a maximal zero-free branch.
Consequently, if

\[
f(x)>0\qquad(x\in I),
\]

then every solution with nonzero initial value exists uniquely on all of \(I\). In particular, if \(I=\mathbb R\), every such solution is entire.
Proof
Suppose, for definiteness, that the branch is defined to the left of \(c\). By continuity, for some \(m>0\),

\[
f(x)\ge m
\]

near \(c\). Hence

\[
u''\ge\frac m u>0.
\]

Thus \(u'\) is increasing.
The continuation criterion gives

\[
\liminf_{x\uparrow c}u(x)=0.
\]

If \(u'\ge0\) somewhere sufficiently near \(c\), monotonicity of \(u'\) would imply that \(u\) is thereafter nondecreasing and bounded away from zero. Thus \(u'<0\) throughout a final interval.
Multiplying

\[
u''\ge\frac m u
\]

by the negative quantity \(u'\) reverses the inequality:

\[
u'u''\le m\frac{u'}u.
\]

Integration from \(x_0\) to \(x\) gives

\[
\frac12\bigl(u'(x)^2-u'(x_0)^2\bigr)
\le
m\log\frac{u(x)}{u(x_0)}.
\]

The right-hand side tends to \(-\infty\) along any sequence for which \(u(x)\to0\), whereas the left-hand side is bounded below by \(-u'(x_0)^2/2\), a contradiction. \(\square\)
Part \(a\) is the special case \(f\equiv1\), for which the additional autonomous first integral makes the solutions explicitly computable.

5.2 Negative limiting forcing: universal collapse law
Theorem 6
Suppose \(c\) is a finite right endpoint of a maximal zero-free branch and

\[
f(c)=-a<0.
\]

Then

\[
|y(x)|\longrightarrow0
\qquad(x\uparrow c),
\]

and

\[
\boxed{
|y(x)|
\sim
(c-x)\sqrt{2a\log\frac1{c-x}}.
}
\]

Moreover,

\[
\boxed{
|y'(x)|
\sim
\sqrt{2a\log\frac1{c-x}}
\longrightarrow+\infty.
}
\]

The analogous statement holds at a left endpoint.
Proof
Put \(u=|y|>0\). Near \(c\),

\[
u''=\frac{f(x)}u<0,
\]

so \(u\) is strictly concave and \(u'\) is decreasing. The continuation criterion gives \(\liminf u=0\). Concavity then implies that \(u\) is eventually strictly decreasing and

\[
u(x)\to0.
\]

Let \(x=x(s)\) be the inverse of \(s=u(x)\) near \(s=0\), and define

\[
p(s)=-u'(x(s))>0.
\]

Since \(ds/dx=u'=-p\),

\[
\frac{dp}{ds}
=
\frac{u''}{p}.
\]

Therefore

\[
\frac{d}{ds}p(s)^2
=2u''
=\frac{2f(x(s))}{s}.
\]

Fix \(s_0>0\). Integration gives

\[
p(s)^2
=
p(s_0)^2
-2\int_s^{s_0}\frac{f(x(r))}{r}\,dr.
\]

Since \(f(x(r))\to-a\),

\[
\boxed{
p(s)^2\sim2a\log\frac1s.
}
\]

Next,

\[
c-x(s)=\int_0^s\frac{dr}{p(r)}.
\]

Using the preceding asymptotic,

\[
c-x(s)
\sim
\frac1{\sqrt{2a}}
\int_0^s\frac{dr}{\sqrt{\log(1/r)}}.
\]

By l'Hôpital's rule,

\[
\int_0^s\frac{dr}{\sqrt{\log(1/r)}}
\sim
\frac{s}{\sqrt{\log(1/s)}}.
\]

Thus

\[
c-x(s)
\sim
\frac{s}{\sqrt{2a\log(1/s)}}.
\]

Taking logarithms shows

\[
\log\frac1s\sim\log\frac1{c-x(s)}.
\]

Inverting the previous relation yields

\[
s
\sim
(c-x(s))
\sqrt{2a\log\frac1{c-x(s)}}.
\]

Since \(s=u(x)=|y(x)|\), this is the asserted formula. The derivative asymptotic follows from the formula for \(p(s)\). \(\square\)

5.3 The case \(f(c)=0\) is genuinely nonuniform
Neither of the previous alternatives extends to \(f(c)=0\).
For example, for \(x<0\), let

\[
y(x)=-x+(-x)^{3/2}.
\]

Then

\[
y''(x)=\frac{3}{4\sqrt{-x}},
\qquad
f(x)=y(x)y''(x)
=\frac34\bigl(\sqrt{-x}-x\bigr).
\]

Defining \(f(0)=0\) gives a continuous forcing, while

\[
y(x)\to0,\qquad y'(x)\to-1,\qquad y''(x)\to+\infty.
\]

Thus the branch cannot extend as a \(C^2\) solution through \(0\), despite \(f(0)=0\) and \(f\ge0\) near \(0\).
This shows that strict positivity in Theorem 5 is necessary.

\begin{enumerate}
\item Exhaustive global structure and failure of uniqueness at flat zeros
\end{enumerate}
Let \(y\in C^2(I)\) be a solution and let

\[
Z=y^{-1}(0).
\]

Then \(Z\) is closed. On every connected component of \(I\setminus Z\),


\(y\) has constant sign;


it is the unique regular branch determined by any one of its nonzero initial data.


Conversely, all classical solutions are obtained by choosing such regular branches on the components of \(I\setminus Z\), setting \(y=0\) on \(Z\subseteq f^{-1}(0)\), and imposing \(C^2\)-matching at the boundary points. The simple and nondegenerate quadratic matching conditions are exactly those in Theorems 3 and 4. Degenerate flat zeros allow additional freedom.
In particular, maximal classical continuation through a zero need not be unique.
Take again

\[
\phi(x)=
\begin{cases}
e^{-1/x^2},&x\ne0,\\
0,&x=0,
\end{cases}
\qquad
f=\phi\phi''.
\]

Define

\[
y_1(x)=\phi(x)
\]

and

\[
y_2(x)=
\begin{cases}
\phi(x),&x\ge0,\\
-\phi(x),&x<0.
\end{cases}
\]

All derivatives of \(\phi\) vanish at \(0\), so both \(y_1,y_2\in C^\infty(\mathbb R)\). For \(x\ne0\),

\[
y_2''(x)=\operatorname{sgn}(x)\phi''(x),
\]

and consequently

\[
y_2y_2''=\phi\phi''=f.
\]

The two solutions agree on \([0,\infty)\) but differ on \((-\infty,0)\). Thus the same regular branch and the same initial data at any \(x_0>0\) admit two distinct entire continuations through the flat zero at \(0\).

\begin{enumerate}
\item Boundary cases
\end{enumerate}
\(f\equiv0\)
Every classical solution of

\[
yy''=0
\]

is affine:

\[
\boxed{y(x)=Ax+B.}
\]

Indeed, \(y''=0\) wherever \(y\ne0\). It also vanishes on the interior of the zero set of \(y\), and continuity of \(y''\) then implies \(y''=0\) at boundary points of that zero set.
\(f\equiv-1\)
This gives the opposite behavior from part \(a\). Define

\[
\Theta(q)=\sqrt2\int_0^q e^{-s^2}\,ds
=\sqrt{\frac{\pi}{2}}\operatorname{erf}(q).
\]

Its range is

\[
\left(-\sqrt{\frac\pi2},\sqrt{\frac\pi2}\right).
\]

All maximal solutions are

\[
x=\tau+a\Theta(q),\qquad
y=\varepsilon a e^{-q^2},
\qquad q\in\mathbb R,
\]

with \(a>0\) and \(\varepsilon=\pm1\). Their maximal domain is the finite interval

\[
\boxed{
\left(
\tau-a\sqrt{\frac\pi2},
\ \tau+a\sqrt{\frac\pi2}
\right).
}
\]

At both endpoints \(y\to0\) and \(|y'|\to\infty\), exactly as predicted by Theorem 6.
Initial data at a zero
For

\[
f(x)=x^2,
\]

both

\[
y_+(x)=\frac{x^2}{\sqrt2},
\qquad
y_-(x)=-\frac{x^2}{\sqrt2}
\]

satisfy the equation and have

\[
y_\pm(0)=y_\pm'(0)=0.
\]

Thus the hypothesis \(y_0\ne0\) in the ordinary uniqueness theorem cannot be omitted.

\begin{enumerate}
\item Final classification
\end{enumerate}
For part \(a\), every maximal solution is entire and is given uniquely by

\[
\boxed{
x=\tau+a\sqrt2\int_0^q e^{s^2}\,ds,
\qquad
y=\varepsilon a e^{q^2},
\quad
a>0,\ \tau\in\mathbb R,\ \varepsilon=\pm1.
}
\]

For part \(b\), under the natural classical interpretation \(y\in C^2\):


\(f\) must be continuous on the solution interval.


Away from zeros of \(y\), the equation has the usual unique maximal zero-free branch for every nonzero initial value.


Zeros can occur only where \(f=0\).


Simple zeros are possible exactly when \(f(x)/(x-c)\) extends continuously.


Nondegenerate tangent zeros are possible exactly when \(f(x)/(x-c)^2\) extends continuously with a positive value; there are then exactly two local tangent solutions.


Degenerate flat zeros may exhibit nonunique continuation.


Between two zeros,

\[
\int_a^b f=-\int_a^b y'^2\le0.
\]



A finite zero-free endpoint with \(f(c)>0\) is impossible.


If \(f(c)<0\), the branch collapses according to

\[
|y(x)|
\sim
|x-c|\sqrt{2|f(c)|\log\frac1{|x-c|}}.
\]



These statements exhaust the universal structure available for an arbitrary prescribed forcing \(f\).
